% !PRINTED 86
\weberpage{86}
%<<CONT
\[
\frac{e_{1,h}i_{1,h}x_1}{i_1} + \frac{e_{2,h}i_{2,h}x_2}{i_2} + \cdots \frac{e_{\nu,h}i_{\nu,h}x_\nu}{i_\nu}
\]
must be integers, or, having regard to the notation (2), with the congruences:
\weq{17}{e_{1,h}i_{1,h}i'_1x_1 + e_{2,h}i_{2,h}i'_2x_2 + \cdots e_{\nu,h}i_{\nu,h}i'_\nu x_\nu \equiv 0\ (\mathrm{mod}\ J).}
These congruences are, as one sees, always satisfied when
$x_1$ is divisible by $i_1$, $x_2$ by $i_2 \dots$, $x_\nu$ by $i_\nu$, and
requirement 2. therefore now reduces to this:

    5. \gesperrt{that the $e_{k,h}$ are to be determined so that the congruences (17) shall be satisfied for no values of the $x$ other than those satisfying the congruences $x_k \equiv 0\ (\mathrm{mod}\ i_k)$.}

    To this is added, as a transformation of requirement 3.:

    6. The numbers
\[
c_{1,h}e_{1,h},\ c_{2,h}e_{2,h},\ \dots\ c_{\nu,h}e_{\nu,h}
\]
\gesperrt{must not all be divisible by $q_h$ when $\varkappa > 1$, and
by $(q_h - 1)$ when $\varkappa = 1$.}

    Requirement 6. contains, on account of 4., no impossibility.

    \gesperrt{By 5. and 6. the necessary and sufficient conditions for the numbers $e_{h,k}$ are expressed.}

    The question that now remains, therefore, is whether, and under what hypotheses, these requirements can be satisfied by a suitable choice of the $e_{h,k}$.
% [sic: Weber switches in these two sentences from the defined e_{k,h} to e_{h,k}; e_{k,h} is intended.]

    This question is greatly simplified by the fact that the
congruence (17) can be split into several. The invariants $i_k$
are, as we have assumed, prime powers. Let one of these
primes be denoted by $p$, and let
\weq{18}{i_1 = p^{\pi_1},\ i_2 = p^{\pi_2},\ \dots\ i_\varrho = p^{\pi_\varrho},}
while the remaining $i_k$, if there are still any, shall no longer
be divisible by $p$. The greatest of the exponents $\pi_1, \pi_2, \dots \pi_\varrho$ we shall denote by $\pi$.

    Then
\[
J = p^\pi J',
\]
and $J'$ is no longer divisible by $p$. Further,
\[
i'_1 = p^{\pi-\pi_1}J',\ i'_2 = p^{\pi-\pi_2}J',\ \dots\ i'_\varrho = p^{\pi-\pi_\varrho}J',
\]
while $i'_{\varrho+1}, \dots i'_\nu$ are divisible by $p^\pi$.
